% pipeline.tex — project pipeline block diagram
% Include via % pipeline.tex — project pipeline block diagram
% Include via % pipeline.tex — project pipeline block diagram
% Include via % pipeline.tex — project pipeline block diagram
% Include via \input{figures/pipeline} from overview.tex

\begin{figure}[H]
\centering
\begin{tikzpicture}[
    >=Stealth,
    block/.style={
        rectangle, draw=black!70, fill=white,
        minimum height=2.2em, minimum width=6.5em,
        text centered, font=\small,
        rounded corners=2pt, thick
    },
    arrow/.style={->, thick, black!70},
    brace label/.style={font=\small\itshape, text=black!70},
    node distance=1.0cm and 0.6cm
]

% --- Pipeline blocks (left to right) ---
\node[block] (corpus) {Corpus};
\node[block, right=of corpus] (kwic) {KWIC extraction};
\node[block, right=of kwic] (mbert) {mBERT attention};
\node[block, right=of mbert] (filt) {Graph filtration};
\node[block, right=of filt] (ph) {\shortstack{Persistent homology\\features}};
\node[block, right=of ph] (comp) {\shortstack{Cross-linguistic\\comparison}};

% --- Arrows ---
\draw[arrow] (corpus) -- (kwic);
\draw[arrow] (kwic) -- (mbert);
\draw[arrow] (mbert) -- (filt);
\draw[arrow] (filt) -- (ph);
\draw[arrow] (ph) -- (comp);

% --- Top brace: linguistic input (corpus + kwic + mbert) ---
\draw[decorate, decoration={brace, amplitude=6pt, raise=8pt}, thick, black!60]
    (corpus.north west) -- (mbert.north east)
    node[midway, above=16pt, brace label] {linguistic input};

% --- Top brace: topology (filt + ph + comp) ---
\draw[decorate, decoration={brace, amplitude=6pt, raise=8pt}, thick, black!60]
    (filt.north west) -- (comp.north east)
    node[midway, above=16pt, brace label] {topology};

\end{tikzpicture}
\caption{%
    The ph-project pipeline. We extract keyword-in-context (KWIC) windows
    from multilingual corpora, run them through mBERT to obtain per-head
    attention matrices, threshold those matrices into a family of graphs
    (the filtration), compute persistent homology features from the
    filtration, and compare the resulting topological signatures across
    languages. The left half is standard NLP; the right half is where
    topological data analysis enters.%
}
\label{fig:pipeline}
\end{figure}
 from overview.tex

\begin{figure}[H]
\centering
\begin{tikzpicture}[
    >=Stealth,
    block/.style={
        rectangle, draw=black!70, fill=white,
        minimum height=2.2em, minimum width=6.5em,
        text centered, font=\small,
        rounded corners=2pt, thick
    },
    arrow/.style={->, thick, black!70},
    brace label/.style={font=\small\itshape, text=black!70},
    node distance=1.0cm and 0.6cm
]

% --- Pipeline blocks (left to right) ---
\node[block] (corpus) {Corpus};
\node[block, right=of corpus] (kwic) {KWIC extraction};
\node[block, right=of kwic] (mbert) {mBERT attention};
\node[block, right=of mbert] (filt) {Graph filtration};
\node[block, right=of filt] (ph) {\shortstack{Persistent homology\\features}};
\node[block, right=of ph] (comp) {\shortstack{Cross-linguistic\\comparison}};

% --- Arrows ---
\draw[arrow] (corpus) -- (kwic);
\draw[arrow] (kwic) -- (mbert);
\draw[arrow] (mbert) -- (filt);
\draw[arrow] (filt) -- (ph);
\draw[arrow] (ph) -- (comp);

% --- Top brace: linguistic input (corpus + kwic + mbert) ---
\draw[decorate, decoration={brace, amplitude=6pt, raise=8pt}, thick, black!60]
    (corpus.north west) -- (mbert.north east)
    node[midway, above=16pt, brace label] {linguistic input};

% --- Top brace: topology (filt + ph + comp) ---
\draw[decorate, decoration={brace, amplitude=6pt, raise=8pt}, thick, black!60]
    (filt.north west) -- (comp.north east)
    node[midway, above=16pt, brace label] {topology};

\end{tikzpicture}
\caption{%
    The ph-project pipeline. We extract keyword-in-context (KWIC) windows
    from multilingual corpora, run them through mBERT to obtain per-head
    attention matrices, threshold those matrices into a family of graphs
    (the filtration), compute persistent homology features from the
    filtration, and compare the resulting topological signatures across
    languages. The left half is standard NLP; the right half is where
    topological data analysis enters.%
}
\label{fig:pipeline}
\end{figure}
 from overview.tex

\begin{figure}[H]
\centering
\begin{tikzpicture}[
    >=Stealth,
    block/.style={
        rectangle, draw=black!70, fill=white,
        minimum height=2.2em, minimum width=6.5em,
        text centered, font=\small,
        rounded corners=2pt, thick
    },
    arrow/.style={->, thick, black!70},
    brace label/.style={font=\small\itshape, text=black!70},
    node distance=1.0cm and 0.6cm
]

% --- Pipeline blocks (left to right) ---
\node[block] (corpus) {Corpus};
\node[block, right=of corpus] (kwic) {KWIC extraction};
\node[block, right=of kwic] (mbert) {mBERT attention};
\node[block, right=of mbert] (filt) {Graph filtration};
\node[block, right=of filt] (ph) {\shortstack{Persistent homology\\features}};
\node[block, right=of ph] (comp) {\shortstack{Cross-linguistic\\comparison}};

% --- Arrows ---
\draw[arrow] (corpus) -- (kwic);
\draw[arrow] (kwic) -- (mbert);
\draw[arrow] (mbert) -- (filt);
\draw[arrow] (filt) -- (ph);
\draw[arrow] (ph) -- (comp);

% --- Top brace: linguistic input (corpus + kwic + mbert) ---
\draw[decorate, decoration={brace, amplitude=6pt, raise=8pt}, thick, black!60]
    (corpus.north west) -- (mbert.north east)
    node[midway, above=16pt, brace label] {linguistic input};

% --- Top brace: topology (filt + ph + comp) ---
\draw[decorate, decoration={brace, amplitude=6pt, raise=8pt}, thick, black!60]
    (filt.north west) -- (comp.north east)
    node[midway, above=16pt, brace label] {topology};

\end{tikzpicture}
\caption{%
    The ph-project pipeline. We extract keyword-in-context (KWIC) windows
    from multilingual corpora, run them through mBERT to obtain per-head
    attention matrices, threshold those matrices into a family of graphs
    (the filtration), compute persistent homology features from the
    filtration, and compare the resulting topological signatures across
    languages. The left half is standard NLP; the right half is where
    topological data analysis enters.%
}
\label{fig:pipeline}
\end{figure}
 from overview.tex

\begin{figure}[H]
\centering
\begin{tikzpicture}[
    >=Stealth,
    block/.style={
        rectangle, draw=black!70, fill=white,
        minimum height=2.2em, minimum width=6.5em,
        text centered, font=\small,
        rounded corners=2pt, thick
    },
    arrow/.style={->, thick, black!70},
    brace label/.style={font=\small\itshape, text=black!70},
    node distance=1.0cm and 0.6cm
]

% --- Pipeline blocks (left to right) ---
\node[block] (corpus) {Corpus};
\node[block, right=of corpus] (kwic) {KWIC extraction};
\node[block, right=of kwic] (mbert) {mBERT attention};
\node[block, right=of mbert] (filt) {Graph filtration};
\node[block, right=of filt] (ph) {\shortstack{Persistent homology\\features}};
\node[block, right=of ph] (comp) {\shortstack{Cross-linguistic\\comparison}};

% --- Arrows ---
\draw[arrow] (corpus) -- (kwic);
\draw[arrow] (kwic) -- (mbert);
\draw[arrow] (mbert) -- (filt);
\draw[arrow] (filt) -- (ph);
\draw[arrow] (ph) -- (comp);

% --- Top brace: linguistic input (corpus + kwic + mbert) ---
\draw[decorate, decoration={brace, amplitude=6pt, raise=8pt}, thick, black!60]
    (corpus.north west) -- (mbert.north east)
    node[midway, above=16pt, brace label] {linguistic input};

% --- Top brace: topology (filt + ph + comp) ---
\draw[decorate, decoration={brace, amplitude=6pt, raise=8pt}, thick, black!60]
    (filt.north west) -- (comp.north east)
    node[midway, above=16pt, brace label] {topology};

\end{tikzpicture}
\caption{%
    The ph-project pipeline. We extract keyword-in-context (KWIC) windows
    from multilingual corpora, run them through mBERT to obtain per-head
    attention matrices, threshold those matrices into a family of graphs
    (the filtration), compute persistent homology features from the
    filtration, and compare the resulting topological signatures across
    languages. The left half is standard NLP; the right half is where
    topological data analysis enters.%
}
\label{fig:pipeline}
\end{figure}
